
SGA 4 XVI 3.7), et,
conjecturalement, si $Y$ et $X$ sont tous deux réguliers (d'après le théorème de pureté cohomologique absolu, démontré tout au moins pour $X$ excellent de caractéristique nulle (SGA 4 XIX 3.2)). Lorsque $X$ et $Y$ sont lisses sur le corps des complexes, on peut également utiliser la classe de Chern entière sur l'espace localement compact $X(\C)$ des points complexes de $X$. La formule (3.1) est établie dans ce cas par \uline{Atiyah} et \uline{Hirzebruch} [1], qui se placent plus généralement dans le cas d'une immersion de variétés analytiques complexes (en travaillant avec le $K^\bullet$ ``naif'') ; il en résulte que la formule (3.1) $\ell$-adique est valable chaque fois que $Y$ et $X$ sont lisses sur un corps \uline{algébriquement clos et de caractéristique nulle}, en utilisant le ``principe de Lefschetz'' et les résultats généraux de SGA 4. Cela rend donc la validité de la formule (3.1) cohomologique extrêmement plausible; mais on notera qu'elle n'est pas démontrée à l'heure actuelle, même si $X$ et $Y$ sont lisses sur un corps $k$ de caractéristique nulle $k$, lorsque $k$ n'est pas algébriquement clos.

\textbf{3.4.} On peut également travailler avec la théorie cohomologique des classes de Chern, style De Rham ou style Hodge [13]. Le problème de la validité de (3.1) dans ce cas ne se pose que si $X$ n'est pas de caractéristique nulle, puisque l'on sait de toutes façons (n\textsuperscript{o}~6) que la formule en question est vraie modulo torsion. Pour les schémas lisses quasi-projectifs sur un corps $k$, la formule (3.1) version De Rham ou Hodge serait d'ailleurs conséquence de la formule (3.1) version classes de cycles mod. équivalence algébrique, de sorte que la version cohomologique De Rham ou Hodge donne une autre façon abordable de tester la validité de (3.1) version ``classes de cycles''.

\textbf{3.5.} La démonstration [1] de la version cohomologique transcendante de (3.1) est de nature essentiellement transcendante, en faisant usage de voisinages tubulaires de $Y$ dans $X$. Cela suggère que la démonstration d'une formule du type (3.1) pourrait être liée à la considération systématique de groupes $K^\bullet$ et de groupes de cohomologie \uline{à support dans} $Y$, considération qui pourrait être un moyen technique suffisant de ``localisation autour de $Y$''.

\section*{4. Relations entre $K^\bullet(X)$ et l'anneau de Chow $A(X)$}

\textbf{4.1.} Supposons pour simplifier que $X$ soit un schéma lisse et quasi-projectif sur un corps $k$ (cf n\textsuperscript{o}~8 pour un cas plus général), de sorte qu'on
